\documentclass[11pt,a4paper]{article}
\usepackage[utf8]{inputenc}
\usepackage{cmap}
\usepackage[T1]{fontenc}
\usepackage[italian]{babel}
\usepackage[margin=2cm]{geometry}
\usepackage[table]{xcolor}
\usepackage{enumitem}
\usepackage{hyperref}
\usepackage{tcolorbox}
\usepackage{titlesec}
\usepackage{booktabs}
\usepackage{tikz}
\usepackage{float}
\usepackage{amsmath}
\usepackage{amssymb}

% Palette di colori professionali (stile TPER)
\definecolor{primaryColor}{HTML}{0054A6}
\definecolor{secondaryColor}{HTML}{4A5568}
\definecolor{backgroundColor}{HTML}{F7FAFC}
\definecolor{borderColor}{HTML}{E2E8F0}
\definecolor{accentGreen}{HTML}{2F855A}
\definecolor{accentOrange}{HTML}{DD6B20}
\definecolor{accentRed}{HTML}{C53030}

\hypersetup{
 colorlinks=true,
 linkcolor=primaryColor,
 urlcolor=primaryColor,
}

\titleformat{\section}
 {\color{primaryColor}\normalfont\Large\bfseries}
 {}{0em}{}
 [\color{borderColor}\titlerule]

\titleformat{\subsection}
 {\color{primaryColor}\normalfont\large\bfseries}
 {}{0em}{}

\titleformat{\subsubsection}[runin]
 {\normalfont\bfseries\color{secondaryColor}}
 {}{0em}{}[:]

\begin{document}

\begin{tcolorbox}[colback=primaryColor,colframe=primaryColor,arc=3mm,halign=center,valign=center,height=2.5cm]
 \color{white}\LARGE\bfseries INFORMATION ARCHITECTURE CARD \\[0.2cm]
 \color{white!80}\large UX Design Project: Help People Help People --- TPER S.p.A. \\[0.1cm]
 \color{white!60}\normalsize Phase C --- Design Proposal
\end{tcolorbox}

\vspace{0.5cm}

% =====================================================================
% Q1: NECESSITÀ DI UNA FASE DI IA SPECIFICA
% =====================================================================
\section{Necessità di una Fase di Information Architecture}

\begin{tcolorbox}[colback=primaryColor!10,colframe=primaryColor,arc=2mm]
\large\textbf{Una fase progettuale specifica per l'architettura dell'informazione dello strumento è stata appropriata e necessaria?} \\[4pt]
\boxed{\text{Sì}} \; \text{Se sì, perché?} \quad \boxed{\text{No}} \; \text{Se no, perché?}
\end{tcolorbox}

\vspace{0.3cm}

\begin{tcolorbox}[colback=backgroundColor,colframe=borderColor,arc=2mm]
{\small
\textbf{Sì, una fase specifica di Information Architecture è stata assolutamente necessaria.} Le evidenze raccolte durante la Phase A (Ethnographic Assessment) e Phase B (Feasibility Study) lo confermano in modo inequivocabile:

\begin{enumerate}[leftmargin=*]
 \item \textbf{Frammentazione dell'ecosistema digitale TPER}: L'ecosistema TPER è composto da canali multipli con funzioni radicalmente diverse, ma questo non è mai reso esplicito all'utente:
 \begin{itemize}
   \item \texttt{tper.it} = sito \textbf{esclusivamente informativo}: pianificatore di viaggio, orari, tipologie di biglietti/abbonamenti, news, contatti, informazioni aziendali. \textbf{Non permette acquisti né transazioni.}
   \item \texttt{solweb.tper.it} = piattaforma dedicata all'acquisto di \textbf{abbonamenti} (solo abbonamenti, non biglietti singoli).
   \item \textbf{Roger app} = app mobile per l'acquisto di \textbf{biglietti singoli} e \textbf{carnet} digitali.
   \item \textbf{Punti Tper} (sportelli fisici) = acquisto di tutti i titoli di viaggio, abbonamenti, assistenza, ritiro tessere.
   \item \textbf{Tabaccherie / PuntoLis} = punti di ricarica per abbonamenti e tessere.
   \item \textbf{Contactless a bordo} (ATM Mastercard/Visa) = pagamento a corsa singola direttamente sul bus.
 \end{itemize}
 Gli utenti non sanno quale canale usare per cosa, e il sito \texttt{tper.it} --- il primo punto di contatto per la maggior parte degli utenti --- non fornisce alcuna guida strutturata a questo ecosistema frammentato.

 \item \textbf{Conferma dai test utente}: Tutti i soggetti testati hanno evidenziato difficoltà strutturali:
 \begin{itemize}
   \item Maria B. (79 anni, competenza minima): ``Non so neanche da dove cominciare'' --- arrivata su \texttt{tper.it}, non ha trovato un modo per acquistare e ha abbandonato.
   \item Sergio M. (67 anni, competenza media): Ha navigato a lungo tra le sezioni di \texttt{tper.it} cercando un pulsante ``Acquista'' --- che su \texttt{tper.it} non esiste --- perdendo oltre 5 minuti.
   \item Chris H. (34 anni, competenza alta): Ha riconosciuto che \texttt{tper.it} è solo informativo, ma ha dovuto indovinare l'esistenza di \texttt{solweb.tper.it} e di Roger per completare i compiti.
   \item Laura P. (70 anni): ``L'interfaccia è troppo ingombrante, c'è troppa confusione'' --- mescola news aziendali, bandi, e informazioni utili senza alcuna gerarchia chiara.
 \end{itemize}

 \item \textbf{Confusione tra informazione e transazione}: Il problema IA centrale è che \texttt{tper.it} mescola contenuti aziendali (news, bandi, appalti, comunicati stampa) con contenuti informativi utili all'utente (orari, tariffe, mappe) senza separazione netta. L'utente che arriva per comprare un biglietto è circondato da informazioni corporate irrilevanti.

 \item \textbf{Impatto diretto sul target HPHP}: Gli utenti fragili (anziani, stranieri, bassa alfabetizzazione digitale) sono i più colpiti. Non solo devono orientarsi in un sito dalle priorità sbagliate, ma devono anche scoprire l'esistenza di canali alternativi (\texttt{solweb}, Roger, Punti Tper) senza alcuna segnaletica che li guidi.

 \item \textbf{Issue D2 -- Violazione G3}: L'analisi euristica ha identificato una violazione grave della linea guida ``Navigazione e architettura informativa''. La navigazione raggiunge fino a quattro livelli di profondità nelle sezioni corporate, mentre le informazioni più cercate (orari, tariffe) sono sepolte sotto etichette aziendali poco chiare.
\end{enumerate}

\textbf{Senza una fase dedicata all'IA, la frammentazione dell'ecosistema --- il problema principale per l'utente --- non potrebbe essere risolta in modo sistematico.}
}
\end{tcolorbox}

\vspace{0.5cm}

% =====================================================================
% Q2: APPROCCIO TOP-DOWN VS BOTTOM-UP
% =====================================================================
\section{Approccio Top-Down vs Bottom-Up}

\begin{tcolorbox}[colback=primaryColor!10,colframe=primaryColor,arc=2mm]
\large\textbf{Se la risposta precedente è stata sì: è stato selezionato un approccio top-down o bottom-up (o entrambi) per l'architettura dell'informazione?}
\end{tcolorbox}

\vspace{0.3cm}

\begin{tcolorbox}[colback=backgroundColor,colframe=borderColor,arc=2mm]
{\small
\renewcommand{\arraystretch}{1.3}
\begin{tabular}{|p{3cm}|p{10cm}|}
\hline
\textbf{Approccio} & \textbf{Selezione} \\
\hline
Top-Down & $\square$ \\
\hline
Bottom-up & $\square$ \\
\hline
Approccio Misto & $\boxtimes$ \textbf{Scelto} \\
\hline
\end{tabular}

\vspace{0.3cm}

\textbf{Giustificazione della scelta:}

L'attuale \texttt{tper.it} adotta un approccio esclusivamente \textbf{top-down istituzionale}: la struttura ricalca l'organigramma aziendale (Azienda, Servizi, Media, Bandi, Contatti) invece di rispondere ai bisogni informativi dell'utente. Il risultato è un'architettura pensata per comunicare l'identità corporate, non per aiutare il cittadino a muoversi.

Il nostro redesign adotta un \textbf{approccio misto}:

\paragraph{Componente Top-Down:} Dagli obiettivi strategici del progetto HPHP discendono le macro-decisioni architetturali:
\begin{itemize}
 \item \textbf{Canale primario = informazione}: \texttt{tper.it} deve essere chiaramente identificato come sito informativo, con una sezione ``Dove Comprare'' che guida l'utente agli altri canali --- senza mai far credere che si possa acquistare direttamente.
 \item \textbf{Guida multicanale strutturata}: L'architettura deve presentare ogni canale dell'ecosistema TPER (\texttt{solweb}, Roger, Punti Tper, tabaccherie, contactless) con vantaggi, svantaggi e requisiti specifici.
 \item \textbf{Supporto multilingua}: Italiano, inglese e arabo per i contenuti informativi principali.
 \item \textbf{Accessibilità e trainer-mode}: Contenuti progettati per essere fruiti da intermediari (operatori Punti Tper, assistenti sociali) che guidano utenti fragili.
 \item \textbf{Pianificazione viaggio come funzione primaria}: Il cuore di \texttt{tper.it} è il route planner --- questo deve essere l'eroe della homepage.
\end{itemize}

\paragraph{Componente Bottom-Up:} Dalle evidenze raccolte nei test utente emergono vincoli precisi:
\begin{itemize}
 \item \textbf{``Dove si compra?'' è la domanda più frequente}: I test dimostrano che la maggior parte degli utenti arriva su \texttt{tper.it} con l'intenzione di acquistare, ma non trova un percorso chiaro. L'architettura deve rispondere a questa domanda \textbf{anche se} l'acquisto non avviene su \texttt{tper.it}.
 \item \textbf{Profondità massima 2 click informativi}: Per ogni informazione su tariffe, orari o tipologie di biglietto, l'utente deve arrivare al contenuto in non più di 2 click.
 \item \textbf{Separazione netta tra informazioni utili e corporate news}: I test hanno dimostrato che mescolare comunicati stampa con orari dei bus disorienta l'utente fragile.
 \item \textbf{Etichette chiare, non aziendali}: ``Abbonamento mensile'' invece di ``Mi Muovo'', ``Dove Comprare'' invece di ``Reti di vendita''.
 \item \textbf{Canale suggerito per ogni esigenza}: Ogni tipologia di titolo di viaggio deve essere accompagnata dall'indicazione esplicita di dove acquistarla (\texttt{solweb}, Roger, Punto Tper, tabaccheria, contactless).
\end{itemize}

\noindent Questo approccio misto è \textbf{radicalmente diverso} dall'attuale: mentre \texttt{tper.it} oggi organizza l'informazione per ``chi comunica il contenuto'', noi la organizziamo per ``cosa l'utente vuole sapere o fare'', bilanciando la completezza informativa con la chiarezza dell'ecosistema multicanale.
}
\end{tcolorbox}

\vspace{0.5cm}

% =====================================================================
% Q3: ORGANIZZAZIONE E LABELING SYSTEM
% =====================================================================
\section{Decisioni su Organizzazione e Sistema di Etichettatura}

\begin{tcolorbox}[colback=primaryColor!10,colframe=primaryColor,arc=2mm]
\large\textbf{Quali sono le principali decisioni relative all'organizzazione e al sistema di etichettatura? Sono diverse dallo strumento esistente?}
\end{tcolorbox}

\vspace{0.3cm}

\begin{tcolorbox}[colback=backgroundColor,colframe=borderColor,arc=2mm]
{\small

Le decisioni sul sistema di organizzazione e labeling sono radicalmente diverse dall'attuale \texttt{tper.it}:

\paragraph{Organizzazione per INTENT dell'utente:}
L'attuale sito organizza i contenuti per ``chi li produce'' (Azienda, Media, Bandi, Servizi). Il redesign organizza i contenuti per ``cosa l'utente vuole fare su tper.it'' (ricordando che tper.it è solo informativo):

\begin{center}
\begin{tcolorbox}[colback=primaryColor!5,colframe=primaryColor,arc=2mm,width=0.95\textwidth]
\small
\textbf{Struttura proposta per tper.it} \\[4pt]
\texttt{[Pianifica]} \;--\; \texttt{[Dove Comprare]} \;--\; \texttt{[Abbonamenti]} \;--\; \texttt{[Biglietti]} \;--\; \texttt{[Aiuto]} \;--\; \texttt{[Azienda]}
\end{tcolorbox}
\end{center}

\begin{enumerate}[leftmargin=*]
 \item \textbf{Pianifica} (priorità \#1): Il cuore di \texttt{tper.it}. Include:
 \begin{itemize}
   \item Pianificatore di viaggio (ricerca percorsi, durata, coincidenze)
   \item Orari e frequenze per linea
   \item Mappe interattive della rete
   \item Stato del servizio in tempo reale (avvisi, ritardi, scioperi)
 \end{itemize}

 \item \textbf{Dove Comprare} (priorità \#2, NOVITÀ ASSOLUTA): Guida completa a tutti i canali dell'ecosistema TPER:
 \begin{itemize}
   \item \texttt{solweb.tper.it}: ``Per abbonamenti online --- vai su solweb.tper.it''
   \item \textbf{Roger app}: ``Per biglietti singoli e carnet digitali --- scarica Roger''
   \item \textbf{Punti Tper}: ``Per tutto: acquisti, abbonamenti, assistenza, ritiro tessere --- trova il punto più vicino''
   \item \textbf{Tabaccherie / PuntoLis}: ``Per ricaricare abbonamenti e tessere''
   \item \textbf{Contactless a bordo}: ``Per pagare la corsa singola con carta o telefono''
   \item \textbf{Tabella comparativa}: Vantaggi, svantaggi, e requisiti di ogni canale
 \end{itemize}

 \item \textbf{Abbonamenti} (priorità \#3): Informazioni complete su:
 \begin{itemize}
   \item Tipologie di abbonamento (mensile, annuale, studenti, over 65, agevolazioni ISEE)
   \item Prezzi e fasce tariffarie
   \item Documenti necessari (ISEE, documento identità, foto)
   \item \textbf{Link esplicito} a \texttt{solweb.tper.it} per l'acquisto online
   \item Guida ``Cosa portare al Punto Tper per fare l'abbonamento''
 \end{itemize}

 \item \textbf{Biglietti} (priorità \#4): Informazioni su:
 \begin{itemize}
   \item Tipologie: biglietto singolo, carnet 10 corse, biglietto giornaliero, biglietto multiplo
   \item Prezzi e zone tariffarie
   \item Dove acquistarli: Roger app, Punti Tper, tabaccherie, contactless a bordo
   \item Validità e condizioni di utilizzo
 \end{itemize}

 \item \textbf{Aiuto} (supporto): FAQ, guide passo-passo per ogni canale, contatti, trainer-mode per intermediari, mappa dei Punti Tper.

 \item \textbf{Azienda} (ultima posizione): Informazioni corporate, news, bandi, appalti, comunicati stampa, contatti istituzionali --- \textbf{separate e dopo} i contenuti utente.
\end{enumerate}

\paragraph{Sistema di etichettatura:}

\begin{enumerate}[leftmargin=*]
 \item \textbf{Mai ``Acquista'' su tper.it}: L'etichetta più importante è ciò che \textbf{non} mettiamo. Nessuna sezione ``Acquista'' o ``Compra'' che farebbe credere all'utente di poter acquistare direttamente. Invece: ``Dove Comprare'' --- etichetta che informa senza promettere transazioni.

 \item \textbf{Da terminologia aziendale a linguaggio utente}:
 \begin{itemize}
   \item ``Pianifica un viaggio'' invece di ``Servizi di mobilità''
   \item ``Dove Comprare'' invece di ``Reti di vendita'' o ``Punti vendita''
   \item ``Abbonamento mensile'' invece di ``Mi Muovo''
   \item ``Hai diritto a uno sconto?'' invece di ``Agevolazioni e delibere regionali''
   \item ``Il tuo abbonamento sta per scadere?'' invece di ``Rinnovo titoli''
 \end{itemize}

 \item \textbf{Supporto multilingua}: Le etichette principali vengono tradotte in inglese e arabo, con priorità ai contenuti di pianificazione e acquisto.

 \item \textbf{Disclosure progressiva}: I termini tecnici (ISEE, DSU, fascia tariffaria, zona) sono accompagnati da spiegazioni inline attivabili, con possibilità di espandere dal riassunto al dettaglio.

 \item \textbf{Segnaletica visiva}: Ogni canale dell'ecosistema è rappresentato da un'icona universale:
 \begin{itemize}
   \item Monitor/globo = \texttt{tper.it} (info)
   \item Schermo = \texttt{solweb.tper.it} (abbonamenti online)
   \item Smartphone = Roger app (biglietti digitali)
   \item Edificio = Punto Tper (sportello fisico)
   \item Negozio = Tabaccheria / PuntoLis (ricariche)
   \item Carta di credito = Contactless a bordo
 \end{itemize}
 Le icone accompagnano ogni riferimento al canale, creando un linguaggio visivo coerente.

\end{enumerate}
}
\end{tcolorbox}

\vspace{0.5cm}

% =====================================================================
% Q4: METADATA FIELDS
% =====================================================================
\section{Campi Metadati Rilevanti}

\begin{tcolorbox}[colback=primaryColor!10,colframe=primaryColor,arc=2mm]
\large\textbf{Quali sono i campi metadati più rilevanti? Sono diversi dallo strumento esistente?}
\end{tcolorbox}

\vspace{0.3cm}

\begin{tcolorbox}[colback=backgroundColor,colframe=borderColor,arc=2mm]
{\small

I metadati sono stati ridefiniti per supportare la navigazione informativa, la guida multicanale, e la ricerca multilingua. Rispetto all'attuale \texttt{tper.it}, che usa metadati pensati per il back-office dei contenuti CMS, il redesign adotta una categorizzazione orientata all'utente e al suo percorso multicanale.

\renewcommand{\arraystretch}{1.2}
\begin{tabular}{|p{3cm}|p{5cm}|p{5cm}|}
\hline
\textbf{Categoria} & \textbf{Metadati} & \textbf{Differenza dall'esistente} \\
\hline
Informazioni di viaggio & \texttt{linea}, \texttt{fermata}, \texttt{zona}, \texttt{durata}, \texttt{orario}, \texttt{frequenza} & Attualmente categorizzati per database interno; noi usiamo attributi per il route planner e la ricerca utente \\
\hline
Canali di acquisto & \texttt{canale} (web/app/fisico/tabaccheria/contactless), \texttt{operazioni}, \texttt{requisiti}, \texttt{orari}, \texttt{link\_esterno} & \textbf{Nuovo}: ogni contenuto su biglietti/abbonamenti è associato a uno o più canali consigliati con link e requisiti \\
\hline
Titoli di viaggio & \texttt{tipo}, \texttt{durata}, \texttt{zona}, \texttt{prezzo}, \texttt{validità}, \texttt{canali\_acquisto}, \texttt{documenti\_necessari} & Nuovo campo \texttt{canali\_acquisto} (array) che elenca dove si può comprare quel titolo; \texttt{documenti\_necessari} per checklist preventiva \\
\hline
Abbonamenti & \texttt{tipo}, \texttt{durata}, \texttt{fasce\_tariffarie}, \texttt{agevolazione}, \texttt{ISEE\_richiesto}, \texttt{canali\_acquisto}, \texttt{link\_solweb} & Nuovo campo \texttt{link\_solweb} per reindirizzamento diretto a \texttt{solweb.tper.it}; campo ISEE per filtro automatico agevolazioni \\
\hline
Guide e tutorial & \texttt{canale\_di\_riferimento}, \texttt{target} (utente/intermediario), \texttt{tempo\_stimato}, \texttt{lingua}, \texttt{passi} & Nuovo: ogni guida è associata al canale che spiega (es. ``Come usare Roger'' $\rightarrow$ canale \texttt{roger}) \\
\hline
Contenuti informativi & \texttt{tipo}, \texttt{target}, \texttt{priorità}, \texttt{data}, \texttt{parole\_chiave}, \texttt{lingua} & Priorizzazione basata sull'intento utente (orari/tariffe prima di news aziendali) \\
\hline
Documenti utente & \texttt{tipo\_documento}, \texttt{scadenza}, \texttt{ente\_rilascio}, \texttt{note} & Metadati per supportare la checklist ``Cosa portare al Punto Tper'' --- mostra all'utente quali documenti servono prima di recarsi allo sportello \\
\hline
\end{tabular}

\vspace{0.3cm}
\textbf{Novità principali rispetto all'esistente:}
\begin{itemize}
 \item \textbf{Metadati ``canale'':} Ogni contenuto su biglietti, abbonamenti o operazioni è etichettato con i canali a cui si riferisce (\texttt{tper.it} info, \texttt{solweb}, Roger, Punti Tper, tabaccherie, contactless). Questo permette alla sezione ``Dove Comprare'' di aggregare dinamicamente tutti i contenuti pertinenti per canale.
 \item \textbf{Metadati di lingua}: Ogni contenuto è associato a versioni linguistiche (IT/EN/AR) per il supporto multilingua.
 \item \textbf{Metadati ``link esterno'':} Per contenuti che rimandano ad altri canali (\texttt{solweb}, download Roger, mappa Punti Tper), il metadato \texttt{link\_esterno} viene usato per generare automaticamente i pulsanti di reindirizzamento.
 \item \textbf{Metadati ``target'':} Distingue contenuti per utente finale da contenuti per intermediario (trainer-mode), consentendo filtri dinamici nell'interfaccia.
\end{itemize}
}
\end{tcolorbox}

\vspace{0.5cm}

% =====================================================================
% Q5: SISTEMA DI NAVIGAZIONE
% =====================================================================
\section{Sistema di Navigazione}

\begin{tcolorbox}[colback=primaryColor!10,colframe=primaryColor,arc=2mm]
\large\textbf{Come pensate di progettare il sistema di navigazione? È diverso dallo strumento esistente?}
\end{tcolorbox}

\vspace{0.3cm}

\begin{tcolorbox}[colback=backgroundColor,colframe=borderColor,arc=2mm]
{\small

Il sistema di navigazione del redesign è profondamente diverso dall'attuale \texttt{tper.it}, che nasconde le funzioni principali dietro etichette aziendali e mescola news con contenuti operativi.

\paragraph{Struttura di navigazione principale:}

\begin{enumerate}[leftmargin=*]

 \item \textbf{Navigazione orizzontale per intento:} Il menu principale è organizzato per azione utente:
 \begin{center}
 \begin{tcolorbox}[colback=primaryColor!5,colframe=primaryColor,arc=2mm,width=0.95\textwidth]
 \small
 \texttt{[Pianifica]} \quad \texttt{[Dove Comprare]} \quad \texttt{[Abbonamenti]} \quad \texttt{[Biglietti]} \quad \texttt{[Aiuto]} \quad $\cdots$ \quad \texttt{[Azienda]}
 \end{tcolorbox}
 \end{center}
 Ogni scheda apre una pagina di primo livello con i contenuti organizzati in schede verticali, senza sottomenu a discesa che aumentano la profondità.

 \item \textbf{Route planner come hero feature}: La homepage di \texttt{tper.it} presenta immediatamente il pianificatore di viaggio come elemento centrale e dominante. Sotto, una sezione ``Stai cercando di comprare?'' con link diretti a \texttt{solweb.tper.it}, Roger, e Punti Tper --- senza mai far credere che si possa acquistare lì.

 \begin{center}
 \begin{tcolorbox}[colback=accentOrange!10,colframe=accentOrange,arc=2mm,width=0.85\textwidth]
 \small \textbf{Banner informativo}: ``\texttt{tper.it} è un sito informativo. Per acquistare abbonamenti vai su \texttt{solweb.tper.it}, per biglietti digitali usa Roger. \\
 \hyperlink{link}{Scopri tutti i canali di acquisto} $\rightarrow$''
 \end{tcolorbox}
 \end{center}

 \item \textbf{Banner di guida canale persistente}: In ogni pagina della sezione ``Abbonamenti'' e ``Biglietti'', un banner superiore informa l'utente su dove può effettivamente acquistare ciò che sta leggendo:
 \begin{itemize}
   \item ``Stai leggendo le informazioni sugli abbonamenti. Per acquistare: vai su solweb.tper.it [pulsante] oppure recati al Punto Tper più vicino [pulsante con mappa]''
   \item ``I biglietti singoli sono disponibili su Roger app [pulsante download] o presso i Punti Tper e tabaccherie [mappa]''
 \end{itemize}

 \item \textbf{Sidebar ``Canali rapidi''}: In ogni pagina informativa su biglietti e abbonamenti, una sidebar fissa mostra i canali disponibili per l'operazione descritta, con icone, tempi stimati e requisiti:
 \begin{center}
 \begin{tcolorbox}[colback=backgroundColor,colframe=borderColor,arc=2mm,width=0.7\textwidth]
 \small
 \begin{tabular}{ll}
 \textbf{Come ottenere un abbonamento}: & \\[4pt]
 $\bullet$ \texttt{solweb.tper.it} & 5 min, carta/tessera sanitaria \\
 $\bullet$ Punto Tper & 10-15 min, documento + ISEE + foto \\
 $\bullet$ Tabaccheria & Solo ricarica, tessera esistente \\
 \end{tabular}
 \end{tcolorbox}
 \end{center}

 \item \textbf{Progress bar per guide passo-passo}: Ogni guida (es. ``Come acquistare il primo abbonamento su solweb'') è accompagnata da una progress bar fissa che mostra i passi:
 \begin{itemize}
   \item Passo 1 di 5: Cosa ti serve (documenti)
   \item Passo 2 di 5: Vai su solweb.tper.it
   \item Passo 3 di 5: Scegli il tipo di abbonamento
   \item Passo 4 di 5: Inserisci i tuoi dati
   \item Passo 5 di 5: Conferma e paga
 \end{itemize}

 \item \textbf{Ponte QR da fisico a digitale}: Ogni Punto Tper espone QR code che l'utente può inquadrare per:
 \begin{itemize}
   \item Aprire la guida passo-passo su \texttt{tper.it} per l'operazione che sta svolgendo
   \item Salvare la lista dei documenti necessari
   \item Avviare una sessione trainer-mode sul proprio telefono
 \end{itemize}

 \item \textbf{Breadcrumb semplificato}: In ogni pagina, un breadcrumb testuale mostra il percorso: \texttt{Home > Abbonamenti > Mensile > Dove acquistare}

\end{enumerate}

\paragraph{Differenze chiave dall'esistente:}

\begin{itemize}
 \item \textbf{Nessuna sezione ``Acquista'':} L'attuale sito ha aree che suggeriscono acquisti impossibili. Il redesign etichetta chiaramente la guida all'acquisto come ``Dove Comprare''.
 \item \textbf{Priorità al route planner}: L'attuale homepage è dominata da news aziendali. La nuova homepage mette il pianificatore di viaggio al centro.
 \item \textbf{Banner di canale vs silos informativi}: L'attuale tratta ogni contenuto come se l'acquisto avvenisse lì. Il redesign aggiunge banner e sidebar che contestualizzano ogni informazione con il canale corretto.
 \item \textbf{Ponte QR}: Nuovo elemento di navigazione che collega il mondo fisico (Punti Tper) a quello digitale (\texttt{tper.it}), assente nell'attuale ecosistema.
 \item \textbf{Profondità contenuti corporate ridotta}: Le informazioni aziendali vengono spostate in fondo alla navigazione, mentre oggi sono in primo piano.
\end{itemize}
}
\end{tcolorbox}

\vspace{0.5cm}

% =====================================================================
% Q6: BROWSING AIDS
% =====================================================================
\section{Ausili alla Navigazione (Browsing Aids)}

\begin{tcolorbox}[colback=primaryColor!10,colframe=primaryColor,arc=2mm]
\large\textbf{Descrivete in breve le vostre scelte riguardo a:}
\end{tcolorbox}

\vspace{0.3cm}

\begin{tcolorbox}[colback=backgroundColor,colframe=borderColor,arc=2mm]
{\small

\subsubsection{Ausili alla navigazione (Browsing aids)}

\begin{enumerate}[leftmargin=*]

 \item \textbf{Tabella comparativa dei canali (NOVITÀ ASSOLUTA)}: In ``Dove Comprare'', una tabella interattiva mostra tutti i canali dell'ecosistema TPER confrontati per:
 \begin{itemize}
   \item Operazioni possibili (biglietti, abbonamenti, ricariche, assistenza)
   \item Requisiti (login, registrazione, documenti fisici)
   \item Tempo stimato
   \item Disponibilità (24/7, orari sportello)
   \item Lingue supportate
   \item Metodi di pagamento accettati
 \end{itemize}
 La tabella è ordinale e filtrabile per esigenza (es. ``Voglio comprare un abbonamento'' $\rightarrow$ mostra solo \texttt{solweb} e Punti Tper).

 \item \textbf{Wizard ``Dove devo andare?'' (NOVITÀ ASSOLUTA)}: Un assistente interattivo a domande che guida l'utente al canale corretto:
 \begin{itemize}
   \item ``Cosa vuoi fare?'' (comprare biglietto / fare abbonamento / ricaricare / assistenza)
   \item ``Come preferisci operare?'' (online / app / dal vivo)
   \item ``Hai già una tessera Tper?'' (sì / no)
   \item ``Hai bisogno di assistenza?'' (sì / no)
 \end{itemize}
 Basandosi sulle risposte, il wizard suggerisce il canale migliore e mostra un percorso personalizzato.

 \item \textbf{Icone visive per ogni canale}: Ogni riferimento a un canale è accompagnato dall'icona corrispondente (monitor, schermo, smartphone, edificio, negozio, carta), creando un linguaggio visivo immediato che aiuta gli utenti con bassa alfabetizzazione o barriera linguistica.

 \item \textbf{Indicatori di posizione}: Ogni pagina mostra chiaramente la sezione corrente:
 \begin{itemize}
   \item Voce di menu attiva evidenziata con colore di sfondo e sottolineatura
   \item Titolo di pagina chiaro e descrittivo
   \item Breadcrumb testuale in alto
   \item Colore di sezione coerente (es. Pianifica = blu, Dove Comprare = verde, Aiuto = arancione)
 \end{itemize}

 \item \textbf{Quick links ``Stai cercando...''}: In ogni pagina della sezione informativa, una box ``Potrebbe interessarti anche'' con link ai contenuti correlati: se stai leggendo degli abbonamenti, potresti voler vedere ``Dove si compra'' o ``Documenti necessari''.

 \item \textbf{Separatori visivi tra contenuto e navigazione}: Uso di tcolorbox, bordi colorati e sfondi differenziati per distinguere chiaramente:
 \begin{itemize}
   \item Contenuto informativo (sfondo bianco)
   \item Guide e wizard (sfondo chiaro con bordo colore canale)
   \item Reindirizzamenti ad altri canali (box evidenziato con icona e pulsante ``Vai a'')
   \item News aziendali (sezione separata, in fondo, con stile differenziato)
 \end{itemize}

\end{enumerate}
}
\end{tcolorbox}

\vspace{0.3cm}

% =====================================================================
% Q7: SEARCH AIDS
% =====================================================================
\section{Ausili alla Ricerca (Search Aids)}

\begin{tcolorbox}[colback=backgroundColor,colframe=borderColor,arc=2mm]
{\small

\subsubsection{Ausili alla ricerca (Search aids)}

\begin{enumerate}[leftmargin=*]

 \item \textbf{Ricerca visibile e task-oriented}: La barra di ricerca è sempre visibile nella header (non nascosta dietro un'icona come nell'attuale). Il placeholder suggerisce azioni: ``Cerca percorsi, orari, canali di acquisto...''

 \item \textbf{Risultati categorizzati per tipo}: I risultati della ricerca sono raggruppati in sezioni visivamente distinte:
 \begin{itemize}
   \item \textbf{Percorsi e orari} (risultati dal route planner)
   \item \textbf{Biglietti e abbonamenti} (informazioni sui titoli di viaggio)
   \item \textbf{Dove comprare} (canali consigliati)
   \item \textbf{Guide e tutorial} (pagine di aiuto passo-passo)
   \item \textbf{News e azienda} (solo in fondo)
 \end{itemize}
 Ogni sezione è collassabile/espandibile e mostra i primi 3 risultati con link ``Vedi tutti''.

 \item \textbf{Autocompletamento con suggerimenti canale}: Mentre l'utente digita, l'autocomplete mostra:
 \begin{itemize}
   \item Suggerimenti di percorso (es. ``da Stazione Centrale a Università'')
   \item Suggerimenti di canale: ``Comprare abbonamento $\rightarrow$ vai su solweb.tper.it''
   \item Suggerimenti di guida: ``Come ricaricare la tessera $\rightarrow$ guide passo-passo''
   \item Suggerimenti di luogo: ``Punto Tper più vicino $\rightarrow$ mappa''
 \end{itemize}

 \item \textbf{Ricerca multilingua con sinonimi}: Il motore accetta query in italiano, inglese e arabo, oltre a sinonimi e termini colloquiali:
 \begin{itemize}
   \item ``sconto'' / ``riduzione'' / ``discount'' / ``agevolazione''
   \item ``abbonamento'' / ``tessera'' / ``pass'' / ``subscription''
   \item ``comprare'' / ``acquistare'' / ``buy'' / ``purchase''
   \item ``dove'' / ``where'' / (arabo: ``ayna'')
 \end{itemize}

 \item \textbf{``Non hai trovato quello che cercavi?''}: In fondo ai risultati, una sezione contestuale che offre:
 \begin{itemize}
   \item Link alla sezione ``Aiuto'' con le FAQ principali
   \item Pulsante ``Avvia wizard Dove Comprare'' per guidare l'utente al canale corretto
   \item Link alla mappa dei Punti Tper per assistenza fisica
   \item Possibilità di contattare l'assistenza telefonica
 \end{itemize}

 \item \textbf{Risultati della ricerca riconoscibili visivamente}: Ogni risultato mostra:
 \begin{itemize}
   \item Icona del canale di riferimento (se pertinente)
   \item Titolo chiaro e descrizione breve
   \item Etichetta di categoria (``Percorso'', ``Guida'', ``Canale'', ``Informazione'')
   \item Link diretto con etichetta d'azione (``Pianifica percorso'', ``Vai a solweb'', ``Scarica Roger'', ``Trova Punto Tper'')
 \end{itemize}

\end{enumerate}
}
\end{tcolorbox}

\vspace{0.3cm}

% =====================================================================
% Q8: CONTENTS AND TASKS
% =====================================================================
\section{Contenuti e Compiti (Contents and Tasks)}

\begin{tcolorbox}[colback=backgroundColor,colframe=borderColor,arc=2mm]
{\small

\subsubsection{Contenuti e compiti (Contents and Tasks)}

I contenuti di \texttt{tper.it} sono organizzati per priorità d'uso, con una gerarchia che riflette i reali bisogni degli utenti emersi dai test:

\begin{enumerate}[leftmargin=*]

 \item \textbf{Priorità \#1 --- Pianificazione del viaggio} (funzione primaria di \texttt{tper.it}):
 \begin{itemize}
   \item Ricerca percorsi: origine $\rightarrow$ destinazione con orari, durata, coincidenze, costo
   \item Orari per linea e fascia oraria
   \item Mappe della rete (urbana, extraurbana, integrata)
   \item Stato del servizio in tempo reale (avvisi, ritardi, scioperi, deviazioni)
   \item Mappa interattiva delle fermate con servizi nelle vicinanze
 \end{itemize}
 \textbf{Accesso}: 1 click dalla homepage. Ricerca testuale con autocomplete delle fermate.

 \item \textbf{Priorità \#2 --- Guida all'acquisto multicanale (Dove Comprare)}:
 \begin{itemize}
   \item Panoramica completa dell'ecosistema TPER con tabella comparativa
   \item \texttt{solweb.tper.it}: guida ``Come acquistare un abbonamento online passo-passo'', requisiti (SPID, carta identità, ISEE), link diretto
   \item \textbf{Roger app}: guida ``Come scaricare e usare Roger'', tipi di biglietti disponibili, metodi di pagamento, link download (App Store / Google Play)
   \item \textbf{Punti Tper}: mappa interattiva con orari, servizi disponibili per ogni punto, tempi di attesa stimati, guida ``Cosa portare''
   \item \textbf{Tabaccherie / PuntoLis}: mappa con indicazione ``Solo ricariche --- non vendono nuovi abbonamenti''
   \item \textbf{Contactless a bordo}: guida ``Come pagare con carta o telefono'', istruzioni per l'uso, limiti e costi
 \end{itemize}

 \item \textbf{Priorità \#3 --- Informazioni su abbonamenti}:
 \begin{itemize}
   \item Tipologie: mensile, annuale, studenti, over 65, agevolazioni ISEE, gratuito per categorie
   \item Prezzi e fasce tariffarie (aggiornati)
   \item Documenti necessari per ogni tipologia
   \item Scadenze e rinnovo
   \item \textbf{Link esplicito} a \texttt{solweb.tper.it} per acquisto online
   \item Guida stampabile ``Cosa portare al Punto Tper''
 \end{itemize}

 \item \textbf{Priorità \#4 --- Informazioni su biglietti}:
 \begin{itemize}
   \item Tipologie: singolo, carnet 10 corse, giornaliero, multiplo, gruppo
   \item Prezzi e zone tariffarie
   \item Validità e condizioni
   \item Dove acquistarli (Roger, Punti Tper, tabaccherie, contactless)
   \item Differenze tra i canali (es. biglietto cartaceo vs digitale)
 \end{itemize}

 \item \textbf{Priorità \#5 --- Aiuto e supporto}:
 \begin{itemize}
   \item FAQ per categoria (abbonamenti, biglietti, ricariche, Roger, contactless)
   \item Guide passo-passo per ogni operazione su ogni canale (trainer-mode)
   \item Glossario dei termini (ISEE, DSU, Fascia, Mi Muovo, Boost, ecc.)
   \item Contatti: telefono, email, PEC, chat (se disponibile)
   \item Mappa interattiva dei Punti Tper e tabaccherie convenzionate
   \item Segnalazione problemi e feedback
 \end{itemize}

 \item \textbf{Ultima priorità --- Contenuti corporate}:
 \begin{itemize}
   \item Chi siamo, mission, valori
   \item News e comunicati stampa
   \item Bandi e appalti
   \item Lavora con noi
   \item Amministrazione trasparente
   \item Privacy e cookie policy
 \end{itemize}
 Questi contenuti sono relegati a una sezione ``Azienda'' separata, con accesso dal footer o da un link secondario nel menu, \textbf{mai} in primo piano.

 \item \textbf{Trainer-mode per intermediari}: Ogni guida ha una modalità ``per operatori'' che aggiunge:
 \begin{itemize}
   \item Note per l'intermediario (es. ``Attenzione: l'utente potrebbe non avere SPID'')
   \item Checklist dei passi da seguire con l'utente
   \item Pulsante ``Stampa riepilogo per l'utente''
   \item Suggerimenti per utenti fragili (es. ``Parla lentamente, indica lo schermo'')
 \end{itemize}
 La trainer-mode è accessibile da un toggle in alto nelle guide, e persiste per l'intera sessione.

\end{enumerate}
}
\end{tcolorbox}

\vspace{0.3cm}

% =====================================================================
% Q9: INVISIBLE COMPONENTS
% =====================================================================
\section{Componenti Invisibili (Invisible Components)}

\begin{tcolorbox}[colback=backgroundColor,colframe=borderColor,arc=2mm]
{\small

\subsubsection{Componenti invisibili}

\begin{enumerate}[leftmargin=*]

 \item \textbf{Persistenza della sessione per trainer-mode}: Il sistema salva automaticamente lo stato delle guide in modalità trainer-mode alla chiusura del browser o al timeout. Al successivo accesso (anche senza login), l'intermediario può riprendere la guida da dove aveva lasciato. Memorizza:
 \begin{itemize}
   \item Passo corrente di ogni guida in corso
   \item Note inserite dall'intermediario
   \item Preferenze di lingua e accessibilità
   \item Documenti segnati come ``verificati'' nella checklist
 \end{itemize}
 La persistenza è realizzata tramite \texttt{localStorage} con fallback su cookie, senza richiedere registrazione. L'intermediario può cancellare la cronologia al termine del percorso assistenziale.

 \item \textbf{Ponte QR code tra Punti Tper e sito web}: Quando un operatore al Punto Tper completa un'assistenza, può generare un QR code stampabile che l'utente porta a casa. Inquadrandolo con lo smartphone, l'utente arriva su \texttt{tper.it} a:
 \begin{itemize}
   \item Un riepilogo di ciò che è stato fatto allo sportello
   \item La guida passo-passo per la prossima operazione (es. ``La prossima volta puoi rinnovare online'')
   \item La lista dei documenti mancanti per la pratica successiva
   \item Un video tutorial sull'uso di \texttt{solweb} o Roger
 \end{itemize}
 Il QR code contiene un identificatore di sessione anonimo (nessun dato personale) valido 30 giorni.

 \item \textbf{Preferenze di lingua persistenti}: Il sistema memorizza la lingua scelta dall'utente (o rilevata dal browser) e la applica a tutte le sessioni successive tramite cookie persistente. Le preferenze sono sincronizzate anche via QR code: se un utente parla arabo e un QR code lo reindirizza a \texttt{tper.it}, il sistema rileva la lingua preferita e mostra la versione localizzata.

 \item \textbf{Analytics dei canali di acquisto (anonimi)}: Il sistema traccia in modo anonimo:
 \begin{itemize}
   \item Quante volte viene visitata la sezione ``Dove Comprare''
   \item Quali canali vengono cliccati più frequentemente
   \item Da quale pagina gli utenti lasciano \texttt{tper.it} per andare su \texttt{solweb.tper.it}
   \item Quali guide sono più consultate
   \item Tasso di completamento delle guide passo-passo
 \end{itemize}
 Questi dati alimentano un dashboard interno (invisibile all'utente) che permette di migliorare continuamente la guida all'acquisto e identificare i canali con maggiore attrito.

 \item \textbf{Generazione automatica del riepilogo stampabile}: Alla fine di ogni guida completata (in particolare in trainer-mode), il sistema genera automaticamente:
 \begin{itemize}
   \item Riepilogo dei passi svolti
   \item Lista dei documenti necessari per l'operazione
   \item Mappa del Punto Tper più vicino (se pertinente)
   \item QR code per accedere al video tutorial
   \item Riferimenti di contatto assistenza TPER
 \end{itemize}
 Il riepilogo è formattato con carattere minimo 14pt, contrasto elevato, e spazio per note scritte a mano.

 \item \textbf{Precaricamento contenuti multilingua}: In background, il sistema precarica le versioni linguistiche dei contenuti più probabili in base alla navigazione corrente, per garantire tempi di risposta rapidi nel passaggio da una lingua all'altra. I contenuti precaricati includono:
 \begin{itemize}
   \item La guida all'acquisto del canale più cliccato nella sessione
   \item Le FAQ della sezione corrente
   \item I contenuti del route planner (universalmente multilingua)
 \end{itemize}

 \item \textbf{Disambiguazione ``Questo sito è informativo''}: Un componente invisibile ma fondamentale è un flag di sessione che, associato a pattern di navigazione (es. utente visita 3+ pagine ``Biglietti'' senza cliccare nessun link di reindirizzamento), attiva un messaggio contestuale:
 \begin{center}
 \begin{tcolorbox}[colback=accentOrange!10,colframe=accentOrange,arc=2mm,width=0.8\textwidth]
 \small \textbf{Stai cercando di acquistare?} Ricorda che su \texttt{tper.it} puoi solo leggere le informazioni. \\[2pt]
 Per acquistare: \texttt{solweb.tper.it} (abbonamenti) $\cdot$ Roger app (biglietti) $\cdot$ Punti Tper (tutto)
 \end{tcolorbox}
 \end{center}
 Questo intervento proattivo riduce l'abbandono per confusione canale, emerso come problema critico nei test.

\end{enumerate}
}
\end{tcolorbox}

\vspace{0.5cm}

% =====================================================================
% RIEPILOGO: KEY DIFFERENCES FROM EXISTING TOOL
% =====================================================================
\section{Riepilogo delle Differenze da tper.it Esistente}

\begin{tcolorbox}[colback=backgroundColor,colframe=primaryColor,arc=2mm]
{\small
\renewcommand{\arraystretch}{1.3}
\begin{tabular}{|p{3.5cm}|p{5cm}|p{5cm}|}
\hline
\textbf{Dimensione} & \textbf{tper.it esistente} & \textbf{Redesign proposto} \\
\hline
\textbf{Natura del sito} & Misto info + transazione (confusione) & \textbf{Esclusivamente informativo}, con guida agli altri canali \\
\hline
\textbf{Homepage} & News aziendali e comunicati stampa & \textbf{Route planner come hero} + banner canali acquisto \\
\hline
\textbf{Organizzazione} & Per struttura aziendale (Azienda, Media, Bandi, Servizi) & Per \textbf{intento utente} (Pianifica, Dove Comprare, Abbonamenti, Biglietti, Aiuto) \\
\hline
\textbf{Etichette} & Terminologia interna (``Mi Muovo'', ``Reti di vendita'', ``Agevolazioni'') & \textbf{Linguaggio utente} + icone per ogni canale \\
\hline
\textbf{Sezione Acquisto} & Nessuna guida strutturata & \textbf{``Dove Comprare''}: tabella comparativa + wizard + guide per canale \\
\hline
\textbf{Ponte canali} & Inesistente (silos separati) & \textbf{Banner, sidebar, QR code} che guidano tra canali \\
\hline
\textbf{Profondità IA} & 4+ livelli navigazione & Massimo 2 livelli per contenuti informativi \\
\hline
\textbf{Contenuti corporate} & In primo piano, mescolati con servizi & Sezione separata ``Azienda'' in \textbf{ultima posizione} \\
\hline
\textbf{Route planner} & Accessibile ma non prioritario & \textbf{Feature principale}, 1 click dalla homepage \\
\hline
\textbf{Multilingua} & Solo italiano & IT/EN/AR per contenuti principali e guide \\
\hline
\textbf{Ricerca} & Nascosta, testuale, senza filtri & Visibile, con \textbf{autocomplete} e risultati \textbf{categorizzati per canale} \\
\hline
\textbf{Wizard canale} & Inesistente & \textbf{``Dove devo andare?''} a domande \\
\hline
\textbf{Ponte QR fisico-digitale} & Inesistente & \textbf{QR code} allo sportello $\rightarrow$ guida online \\
\hline
\textbf{Trainer-mode} & Inesistente & Guide con note intermediario, checklist, riepilogo stampabile \\
\hline
\textbf{Sessioni} & Stateless & Persistenza trainer-mode + analisi anonima dei canali \\
\hline
\end{tabular}
}
\end{tcolorbox}

\vspace{0.5cm}

\begin{center}
{\color{secondaryColor}\small Documento compilato sulla base delle evidenze raccolte durante la Phase A (Ethnographic Assessment) e Phase B (Feasibility Study) del progetto \\[4pt]
\textbf{Phase C -- Design Proposal: Information Architecture Card} \\[4pt]
Progetto: UX Design Project: Help People Help People -- TPER S.p.A.}

\vspace{0.3cm}
{\color{secondaryColor}\itshape\small
\texttt{tper.it} = sito informativo \;$\cdot$\; \texttt{solweb.tper.it} = abbonamenti online \;$\cdot$\; Roger app = biglietti digitali \;$\cdot$\; Punti Tper = tutto \;$\cdot$\; Tabaccherie = ricariche \;$\cdot$\; Contactless = corsa singola
}
\end{center}

\end{document}
